\section{Algebraic categorical models}\label{sec:type-theoretic-axioms}

In this section,
we introduce axioms for a model of directed type theory
in the style of algebraic type theory,
that could be satisfied in models other than
the one in $\Cat$ considered so far (\Cref{ex:ACM-of-opfibrations}).
Once the axioms have all been stated,
we will check that they are indeed satisfied in $\Cat$.

The algebraic style allows us to state our axioms in a concise manner
that automatically guarantees stability under substitution of each construction,
so that when a syntax is fully worked out,
one can easily define an interpretation.
We begin with an axiom for the
external fibration structure on a category.

\section*{Opfibrations and fibrations} 
% \label{sec:types-opfibration-and-fibration}
\begin{axiom}\label{ax:opfibration-preclan}
$\CC$ is a large category with a terminal object and a
class of maps $\opfib$ forming a preclan $(\CC,\opfib)$.
We call maps in $\opfib$ \emph{opfibrations}\footnote{
  Since we are now working in an axiomatic setting,
  $\opfib$ could consist of something other than the class of split opfibrations in $\Cat$,
  and so we drop the adjective ``split''.
  Indeed $\opfib$ could be interpreted in $\Cat$
  as the class of cloven opfibrations
  or discrete opfibrations.
  It could also be interpreted in $\Grpd$
  as the class of split or cloven isofibrations.
}.
\end{axiom}

As usual, objects in the category $\CC$ are regarded as contexts,
and morphisms as substitutions.

\begin{axiom}\label{ax:global-op}
  We have an endofunctor $\Op{} : \CC \to \CC$
  satisfying $\Op \circ \Op = \id_\CC$.
\end{axiom}

The functor $\Op : \CC \to \CC$ is self-inverse,
hence preserves all limits, including the terminal object.
\[ \Op{1} \iso 1\]

\begin{definition}
  We write $\fib$ for the class of maps
  $f : X \to Y$ such that $\Op{f} : \Op{X} \to \Op{Y}$
  is an opfibration.
  Maps in $\fib$ are called \emph{fibrations}.
\end{definition}

It follows from 
\Cref{ax:opfibration-preclan} and \Cref{ax:global-op}
that $(\CC,\fib)$ is also a preclan.

\begin{proposition}\label{prop:op-preclan-morphism}
  The functor $\Op : \CC \to \CC$
  extends to two variance-reversing preclan morphisms
  \[ \Op{} : (\CC, \opfib) \to (\CC, \fib)
  \quad \quad
  \Op{} : (\CC, \fib) \to (\CC, \opfib) \]
  which are inverses of each other.
\end{proposition}

Recall the definitions of preclan-exponentiability
(\Cref{def:exponentiable})
and $\tau$-preclans (\Cref{def:double-preclan}).
The following proposition says that given half of the exponentiability conditions we require,
$\Op$ automatically generates the other half.

\begin{proposition}\label{prop:exponentiability-opposite}
  The following conditions are equivalent
  \begin{enumerate}
  \item Fibrations are opfibration-exponentiable, meaning
    \[ \fib \subseteq \Exp_\opfib \]
  \item Opfibrations are fibration-exponentiable.
    \[\opfib \subseteq \Exp_\fib\]
  \item Opfibrations $(\CC,\opfib)$ and fibrations $(\CC,\fib)$ form a $\tau$-preclan.
  \end{enumerate}
\end{proposition}
\begin{proof}
  (1) $\Rightarrow$ (2).
  Let $f : Y \to X$ be an opfibration.
  By \Cref{lem:stable-class-of-exponentiable-maps}
  it suffices to show that the pullback functor $f^* : \CC / X \to \CC / Y$
  has a partial right adjoint $f_* : \fib(Y) \to \fib(X)$.

  The self-inverse $\Op : \CC \to \CC$ induces isomorphisms between slices
  \[\Op / Y : \CC / Y \iso \CC / \Op{Y} 
  \quad \text{ and } \quad
  \Op / X : \CC / X \iso \CC / \Op{X} \]
  Since $\Op : \CC \to \CC$ preserves pullbacks,
  $f^* : \CC / X \to \CC / Y$
  is isomorphic under the above equivalences to
  $(\Op{f})^* : \CC / \Op{X} \to \CC / \Op{Y}$.
  \[
  \begin{tikzcd}
    {\CC / {Y}} & {\CC / \Op{Y}} & {\opfib(\Op{Y})} & {\fib(Y)} \\
    {\CC / {X}} & {\CC / \Op{X}} & {\opfib(\Op{X})} & {\fib(X)}
    \arrow["\iso"{description}, draw=none, from=1-1, to=1-2]
    \arrow[hook', from=1-3, to=1-2]
    \arrow["\iso"{description}, draw=none, from=1-3, to=1-4]
    \arrow[""{name=0, anchor=center, inner sep=0}, "{(\Op{f})_*}", from=1-3, to=2-3]
    \arrow["{f_*}", dashed, from=1-4, to=2-4]
    \arrow["{f^*}", from=2-1, to=1-1]
    \arrow["\iso"{description}, draw=none, from=2-1, to=2-2]
    \arrow[""{name=1, anchor=center, inner sep=0}, "{(\Op{f})^*}", from=2-2, to=1-2]
    \arrow[hook', from=2-3, to=2-2]
    \arrow["\iso"{description}, draw=none, from=2-3, to=2-4]
    \arrow["{\dashv_\partial}"{description}, draw=none, from=1, to=0]
  \end{tikzcd}
  \]
  The preclan isomorphism
  $\Op : (\CC, \fib) \to (\CC, \opfib)$
  also induces isomorphisms
  \[
  \Op{} / Y : \fib(Y) \iso \opfib(\Op{Y})
  \quad \text{ and } \quad
  \Op{} / X : \fib(X) \iso \opfib(\Op{X})
  \]
  Providing a partial right adjoint
  $f_* : \fib(Y) \to \fib(X)$ of $f^* : \CC / X \to \CC / Y$
  amounts to providing a partial right adjoint
  $(\Op{f})_*$ of $(\Op{f})^*$, which we have by (1).
  
  (2) $\Rightarrow$ (1) is similar.
  (1) $\Leftrightarrow$ (3) immediately follows from (1) $\Rightarrow$ (2).
\end{proof}

\begin{axiom}\label{ax:exponentiability}
  The equivalent conditions of \Cref{prop:exponentiability-opposite} hold.
\end{axiom}

% \begin{axiom}
% $\CC$ is a category with a terminal object with two preclans
% $(\CC,\opfib)$ and $(\CC,\fib)$
% that are exponentiable against each other (see \Cref{def:exponentiable}):
% \[\fib \subseteq \Exp_\opfib
%   \quad \text{ and } \quad
%   \opfib \subseteq \Exp_\fib\]
% \end{axiom}

\section*{Vertical opposites}

In the $\Cat$ model (\Cref{ex:ACM-of-opfibrations}), we need to be able to place restrictions
on the morphisms between opfibrations -- they should be
opcartesian functors.
Let us assume the axioms that have been introduced so far are available to us.
Consider the pseudofunctor $\opfib(-) : \CC^\op \to \tCat$,
that takes an object $X$ in $\CC$ to the category $\opfib(X)$
of opfibrant objects in the slice $\CC / X$
and a morphism $f : X \to Y$ to the pullback functor
$f^* : \opfib(Y) \to \opfib(X)$.
% Similarly, $\fib(-) : \CC^\op \to \Cat$.

\begin{axiom}\label{ax:opcartesian-map}
  For each $X \in \CC$,
  there is a wide subcategory $\opfibcart(X) \to \opfib(X)$,
  forming a pseudofunctor $\opfibcart(-) : \CC^\op \to \tCat$.
  We call the maps in $\opfibcart(X)$ \emph{opcartesian maps} over $X$.
  Furthermore, all maps between opfibrations over $1$ are opcartesian.
  \[\opfibcart(1) = \opfib(1)\]
\end{axiom}

\begin{axiom}\label{ax:local-op}
  The following structure is present in $\CC$.
  \begin{itemize}
    \item We have a pseudotransformation $\vOp{} : \opfibcart(-) \to \opfibcart(-)$ 
    \item For each object $X$ in $\CC$,
    $\vOp[X]{} : \opfibcart(X) \to \opfibcart(X)$ is self inverse.
    \[\vOp[X]{} \circ \vOp[X]{} \iso \id_{\opfibcart(X)}\]
    \item The functor $\vOp[1]{} : \opfibcart(1) \to \opfibcart(1)$
    % and $\vOp[1]{} : \fibcart(1) \to \fibcart(1)$ are both (isomorphic to)
    is (up to isomorphism) a restriction of the global opposite functor
    $\Op{} : \CC \to \CC$.
    \[ \begin{tikzcd}
        {\opfibcart(1)} & {\opfibcart(1)} \\
        \CC & \CC
        \arrow["{{{\vOp[1]{}}}}", from=1-1, to=1-2]
        \arrow[hook, from=1-1, to=2-1]
        \arrow[hook, from=1-2, to=2-2]
        \arrow["{\Op{}}"', from=2-1, to=2-2]
      \end{tikzcd} \]
    % \[\begin{tikzcd}
    %     {\opfibcart (1)} & {\opfibcart (1)} \\
    %     \CC & \CC \\
    %     {\fibcart(1)} & {\fibcart(1)}
    %     \arrow["{{\vOp[1]{}}}", from=1-1, to=1-2]
    %     \arrow[hook, from=1-1, to=2-1]
    %     \arrow[hook, from=1-2, to=2-2]
    %     \arrow["\Op{}", from=2-1, to=2-2]
    %     \arrow[hook, from=3-1, to=2-1]
    %     \arrow["{{\vOp[1]{}}}"', from=3-1, to=3-2]
    %     \arrow[hook, from=3-2, to=2-2]
    %   \end{tikzcd}\]
  \end{itemize}
\end{axiom}

For an opfibration $A \to X$, the opfibration $\vOp[X]{A} \to X$
is the \emph{vertical opposite} of $A$.
In $\Cat$, the vertical opposite of $A$ can be thought of as
the category where,
over each point $x$ in $X$,
paths in each fibre $A_x$ are reversed.
This is discussed in detail in \Cref{sec:preliminaries}.

We construct the duals to \Cref{ax:opcartesian-map} and \Cref{ax:local-op}.
\begin{definition}
  Define the pseudofunctor $\fibcart(-) : \CC^\op \to \tCat$ such that
  for each $X$ in $\CC$, $\fibcart(X)$ is the wide subcategory of $\fib(X)$ consisting of maps
  $F : A \to B$ over $X$ such that $\Op F : \Op A \to \Op B$ is
  an opcartesian in $\opfibcart(\Op X)$.
  We call maps in $\fibcart(X)$ \emph{cartesian maps} over $X$.

  Define the pseudotransformation $\vOp{} : \fibcart(-) \to \fibcart(-)$
  such that the component for each $X$ in $\CC$ is the composition
  \[ \begin{tikzcd}[column sep = large]
    {\fibcart(X)} & {\opfibcart (\Op X)} & {\opfibcart (\Op X)} & {\fibcart(X)}
    \arrow["{\Op / X}", from=1-1, to=1-2]
    \arrow["\sim"', draw=none, from=1-1, to=1-2]
    \arrow["{\vOp[\Op X]{}}", from=1-2, to=1-3]
    \arrow["{\Op / \Op X}", from=1-3, to=1-4]
    \arrow["\sim"', from=1-3, to=1-4]
  \end{tikzcd} \]
\end{definition}

\begin{proposition}
  For each object $X$ in $\CC$,
  \[\vOp[X]{} \circ \vOp[X]{} \iso \id_{\fibcart(X)} : \fibcart(X) \to \fibcart(X)\]
    % \item These operations agree up to isomorphism where their domains intersect.
    % \[
    %   \begin{tikzcd}
    %   {\opfibcart (X)} & {\opfibcart (X)} \\
    %   {\opfibcart  \cap \fib (X)} & {\opfibcart  \cap \fib(X)} \\
    %   {\fib(X)} & {\fib(X)}
    %   \arrow["{\vOp[X]{}}", from=1-1, to=1-2]
    %   \arrow[hook, from=2-1, to=1-1]
    %   \arrow[dashed, from=2-1, to=2-2]
    %   \arrow[hook, from=2-1, to=3-1]
    %   \arrow[hook, from=2-2, to=1-2]
    %   \arrow[hook, from=2-2, to=3-2]
    %   \arrow["{\vOp[X]{}}"', from=3-1, to=3-2]
    % \end{tikzcd}
    % \]
    and $\vOp[1]{} : \fibcart(1) \to \fibcart(1)$
    is (up to isomorphism) a restriction of the global opposite functor
    $\Op{} : \CC \to \CC$.
    \[ \begin{tikzcd}
        {\fibcart(1)} & {\fibcart(1)} \\
        \CC & \CC
        \arrow["{{{\vOp[1]{}}}}", from=1-1, to=1-2]
        \arrow[hook, from=1-1, to=2-1]
        \arrow[hook, from=1-2, to=2-2]
        \arrow["{\Op{}}"', from=2-1, to=2-2]
      \end{tikzcd} \]
\end{proposition}
\begin{proof}
  For the second isomorphism, consider the following diagram.
  \[ \begin{tikzcd}
    A & {\Op {A}} && {\Op {A}} & A \\
    {\fibcart(1)} & {\opfibcart (\Op{1})} & {\opfibcart (1)} & \CC & \CC \\
    {\fibcart(1)} & {\opfibcart (\Op{1})} & {\opfibcart (1)} & \CC & \CC
    \arrow[maps to, from=1-1, to=1-2]
    \arrow[maps to, from=1-2, to=1-4]
    \arrow[maps to, from=1-4, to=1-5]
    \arrow["\sim", from=2-1, to=2-2]
    \arrow["{\vOp[1]{}}"', from=2-1, to=3-1]
    \arrow["\sim", from=2-2, to=2-3]
    \arrow["{\vOp[\Op{1}]{}}", from=2-2, to=3-2]
    \arrow[hook, from=2-3, to=2-4]
    \arrow["{\vOp[1]{}}", from=2-3, to=3-3]
    \arrow["\Op{}", from=2-4, to=2-5]
    \arrow["{\Op{}}", from=2-4, to=3-4]
    \arrow["\Op{}", from=2-5, to=3-5]
    \arrow["\sim"', from=3-1, to=3-2]
    \arrow["\sim"', from=3-2, to=3-3]
    \arrow[hook, from=3-3, to=3-4]
    \arrow["\Op{}"', from=3-4, to=3-5]
  \end{tikzcd}\]
  The first square is the definition of $\vOp[] : \fibcart(-) \to \fibcart(-)$;
  the second is pseudonaturality of $\vOp[] : \opfibcart (-) \to \opfibcart (-)$
  applied to the isomorphism $1 \iso \Op 1$;
  the third square is from \Cref{ax:local-op}.
  A straightforward computation shows that the composition of the horizontal maps 
  is the inclusion $\fibcart(1) \hookrightarrow \CC$.
\end{proof}

\begin{proposition}
  Opfibrations over $1$ are the same as fibrations over $1$.
  All maps between fibrations over $1$ are cartesian.
  \[ \opfib(1) = \opfibcart(1) = \fibcart(1) = \fib(1) \]
\end{proposition}
\begin{proof}
  Let $!_X : X \to 1$ be an opfibration.
  By \Cref{ax:local-op}, viewing $X$ as an object in $\opfibcart(1)$,
  we have $\Op{X} \iso \vOp[1]{X}$ is in $\opfibcart(1)$,
  making the unique map $\Op{X} \to 1$ an opfibration.
  It follows from \Cref{prop:op-preclan-morphism} that the unique map
  $\Op{\Op{X}} \to \Op{1}$ is a fibration.
  Since $\Op{} \circ \Op{} = \id_\CC$,
  this says that $X \to 1$ is a fibration.
\end{proof}

We will say that an object $X$ in $\CC$ is \emph{fibrant} if it belongs
to $\opfibcart(1)$.

\begin{definition}\label{def:ACM}
  An \emph{op structure} (for an algebraic categorical model) consists of the structures
  \[(\CC,\opfibcart,\Op,\vOp)\]
  satisfying the axioms
  \labelcref{ax:opfibration-preclan,ax:global-op,ax:exponentiability,,ax:opcartesian-map,ax:local-op}.
\end{definition}

\section*{Universes}
From now on, we fix an op structure $(\CC,\opfibcart,\Op,\vOp)$.
% Like in \Cref{def:elementary-universe},
\begin{definition}
We say that a morphism $u : U' \to U$ in $\CC$
is a \emph{universe} if for any map $A : \Gamma \to U$
there is a chosen pullback of $u$ along $A$
(cf. \Cref{def:elementary-universe}).
\end{definition}

\begin{axiom}\label{ax:universe}
  We have a universe $\ulow : \Ulow \to \U$
  such that $\U \to 1$ is fibrant,
  and $\ulow : \Ulow \to \U$ is an opfibration.
  We call $\ulow : \Ulow \to \U$ the \emph{universe of small opfibrations}.
\end{axiom}

Since $\U$ is fibrant (hence $\U \in \opfib(1) = \fib(1)$) and $\ulow$ is an opfibration,
$\Ulow$ is also fibrant.
Any opfibration is $\fib$-exponentiable
and $\U$ is fibrant,
so $(\ulow : \Ulow \to \U)$ is also a polynomial signature in the preclan
$(\CC, \fib)$.
\[
\begin{tikzcd}
	{\opfib(1) \ni} & \Ulow & \in \fib(1) \\
	{\opfib \ni} && {\in \Exp_\fib} \\
	{\opfib(1) \ni} & \U & {\in \fib(1)}
	\arrow["\ulow", from=1-2, to=3-2]
\end{tikzcd}
\]

For a map $A : \Gamma \to \U$ in $\CC$,
we denote the chosen pullback of $\ulow : \Ulow \to \U$ along
$A$ as $d_A : \Gamma \opext A \to \Gamma$, and refer to it as the 
\emph{opfibrant context extension}, or \emph{opextension} for short.
\[ \begin{tikzcd}
    {\Gamma \opext A} & \Ulow \\
    \Gamma & \U
    \arrow[from=1-1, to=1-2]
    \arrow["{d_A}"', from=1-1, to=2-1]
    \arrow["\lrcorner"{anchor=center, pos=0.125}, draw=none, from=1-1, to=2-2]
    \arrow["\ulow", from=1-2, to=2-2]
    \arrow["A"', from=2-1, to=2-2]
  \end{tikzcd}\]

\begin{definition}
  We write $\vupp : \Vupp \to \V$ for
  the image of $\ulow : \Ulow \to \U$ under the functor
  $\Op : \CC \to \CC$.
  \[
    \V := \Op{\U} \quad
    \Vupp := \Op{\Ulow} \quad
    \vupp := \Op{\ulow}
  \]
  % \[\vlow : \Vlow \to \V \quad \quad := \quad \quad
  % \Op{\ulow} : \Op{\Ulow} \to \Op{\U}\]
\end{definition}

It follows from \Cref{ax:global-op} and \Cref{ax:local-op} that
$\V$ is fibrant and $\vupp : \Vupp \to \V$ is an fibration.

As we will see in \Cref{ex:groupoidal-opfibrations,ex:discrete-opfibrations},
we might want several different universes at the same time,
and not just stacked up in a hierarchy.
To keep things simple, however,
we will only work with a single one 
(along with its various opposites, of course).

\section*{Small opfibrations and small fibrations}
Now that we are equipped with a universe of small opfibrations
$\ulow : \Ulow \to \U$,
we can consider the class of maps classified by the universe.
Consider the subclass of maps $\opfib_\U \subseteq \opfib$
consisting of those maps that are pullbacks of $\ulow$
along some map into $\U$.
Naturally, we call the maps in $\opfib_\U$ \emph{small opfibrations}.
\[
  \opfib_\U := \{ f : E \to B \st \begin{tikzcd}
	E & \Ulow \\
	B & \U
	\arrow["\exists", dashed, from=1-1, to=1-2]
	\arrow["f"', from=1-1, to=2-1]
	\arrow["\lrcorner"{anchor=center, pos=0.125}, draw=none, from=1-1, to=2-2]
	\arrow["\ulow", from=1-2, to=2-2]
	\arrow["\exists"', dashed, from=2-1, to=2-2]
\end{tikzcd}\}
\]
Dually, we have a class of \emph{small fibrations} $\fib_\V$,
which is the class of pullbacks of $\vupp : \Vupp \to \V$.

\begin{proposition}\label{prop:op-preclan-morphism-small}
  A map $f : E \to B$ is a small fibration if and only if
  $\Op{f} : \Op{E} \to \Op{B}$ is a small opfibration.
  Hence we have two preclan morphisms
  \[ \Op{} : (\CC, \opfib_\U) \to (\CC, \fib_\V)
  \quad \quad
  \Op{} : (\CC, \fib_\V) \to (\CC, \opfib_\U) \]
\end{proposition}
\begin{proof}
  This follows from the fact that $\Op$ preserves pullbacks.
\end{proof}

% The following axiom is trivial if one assumes the axiom of choice.
% when picking a pullback square when classifying a small (op)fibration.
In the $\Cat$ model (\Cref{ex:ACM-of-opfibrations}),
we could interpret $\opfib$ as the class of split opfibrations.
Unfortunately,
since we are treating $\opfib$ as a class of maps,
i.e. maps satisfying a property,
rather than maps equipped with structure, 
we will need to use the axiom of choice
to pick classifying maps of split opfibrations.
However, a careful setup of \Cref{sec:exponentiable}
and \Cref{sec:polynomials}
that replaces the class of maps $\opfib$ with a
(not necessarily full or faithful)
functor $\opfib \to \CC^\two$ should be possible.
We leave this technical solution to future work.
For now, let us use the following axiom to keep track of when we might need
the axiom of choice.

\begin{axiom}\label{ax:classifying-map}
  Every small opfibration $A : X \to \Gamma$ in $\opfib_\U$
  is equipped with a chosen \emph{classifying map}
  $\ulcorner A \urcorner : \Gamma \to \U$,
  so that $A : X \to \Gamma$ is isomorphic to the opextension
  $\Gamma \opext \ulcorner A \urcorner \to \Gamma$ in $\opfib(\Gamma)$.
  \[ \begin{tikzcd}
      X & {\Gamma \opext \ulcorner A \urcorner} & \Ulow \\
      & \Gamma & \U
      \arrow["\sim", from=1-1, to=1-2]
      \arrow["A"', from=1-1, to=2-2]
      \arrow[from=1-2, to=1-3]
      \arrow[from=1-2, to=2-2]
      \arrow["\lrcorner"{anchor=center, pos=0.125}, draw=none, from=1-2, to=2-3]
      \arrow["\ulow", from=1-3, to=2-3]
      \arrow["{\ulcorner A \urcorner}"', from=2-2, to=2-3]
    \end{tikzcd} \]
\end{axiom}

\section*{Vertically opposite universes}

% It is convenient to have notation for the vertical opposites of the two universes
% $\ulow : \Ulow \to \U$ and $\vlow : \Vlow \to \V$.
Viewing $\ulow : \Ulow \to \U$ as an object in $\opfibcart(\U)$,
we can apply $\vOp[\U]{} : \opfibcart(\U) \to \opfibcart(\U)$ to it,
obtaining an opfibration
\[\uupp : \Uupp \to \U\]
Dually, applying $\vOp[\V]{} : \fibcart(\V) \to \fibcart(\V)$ to the object
$\vupp : \Vupp \to \V$,
we have a fibration
% Viewing $\vlow : \Vlow \to \V$ as an object in $\fibcart(\V)$,
% we can apply  to it,
% obtaining a fibration 
\[\vlow : \Vlow \to \V\]

\begin{lemma}\label{lem:vupp-eq}
  $\Vlow = \Op{\Uupp}$ and $\vlow = \Op{\uupp} : \Vlow \to \V$.
\end{lemma}
\begin{proof}
  \[ \vlow = \vOp[\V]{\vupp}
  = \vOp[\V]{\Op \ulow}
  = \Op (\vOp[\U]{\Op \Op \ulow})
  = \Op (\vOp[\U]{\ulow})
  = \Op {\uupp}\]
\end{proof}

We write $\opfibcart_\U(X)$ for the wide subcategory of $\opfib_\U(X)$
consisting only of opcartesian maps.
\[ \begin{tikzcd}
	{\opfibcart_\U (X)} & {\opfibcart (X)} \\
	{\opfib_\U (X)} & {\opfib (X)}
	\arrow[hook, from=1-1, to=1-2]
	\arrow[from=1-1, to=2-1]
	\arrow["\lrcorner"{anchor=center, pos=0.125}, draw=none, from=1-1, to=2-2]
	\arrow[from=1-2, to=2-2]
	\arrow[hook, from=2-1, to=2-2]
\end{tikzcd} \]

\begin{proposition}\label{prop:vertical-opposite-small}
  The following are equivalent
  \begin{enumerate}
    \item $\vOp[X] : \opfibcart(X) \to \opfibcart(X)$ preserves $\opfibcart_\U(X)$,
    for every $X$ in $\CC$.
    \[ \begin{tikzcd}
      {\opfibcart_\U(X)} & {\opfibcart_\U(X)} \\
      {\opfibcart(X)} & {\opfibcart(X)}
      \arrow[dashed, from=1-1, to=1-2]
      \arrow[hook, from=1-1, to=2-1]
      \arrow[hook, from=1-2, to=2-2]
      \arrow["{\vOp[X]{}}"', from=2-1, to=2-2]
    \end{tikzcd} \]
    \item $\vOp[\U] : \opfibcart(\U) \to \opfibcart(\U)$ preserves $\opfibcart_\U(\U)$.
    \item $\uupp : \Uupp \to \U$ is a small opfibration.
    \item Each $\vOp[X] : \fibcart(X) \to \fibcart(X)$ preserves $\fibcart_\V(X)$.
    \item $\vOp[\V] : \fibcart(\V) \to \fibcart(\V)$ preserves $\fibcart_\V(\V)$.
    \item $\vupp : \Vupp \to \V$ is a small fibration.
  \end{enumerate}
\end{proposition}
\begin{proof}
  (1) $\Rightarrow$ (2) $\Rightarrow$ (3) is clear.
  We prove (3) $\Rightarrow$ (1)
  Suppose $A$ is an object in $\opfibcart(X)$.
  By pseudonaturality of $\vOp[X]{}$,
  \[ \vOp[X]{A} \iso \vOp[X]{(\ulcorner A \urcorner)^* \ulow}
    \iso (\ulcorner A \urcorner)^* \vOp[\U]{\ulow}
    = (\ulcorner A \urcorner)^* \uupp \]
  Hence $\vOp[X]{A}$ is a pullback of a small opfibration,
  which is small.

  % So far we have shown that $1. \Leftrightarrow 2. \Leftrightarrow 3.$
  By duality, we also have (4) $\Leftrightarrow$ (5) $\Leftrightarrow$ (6).
  It remains to check that (3) $\Leftrightarrow$ (6).
  Indeed, this follows from 
  \Cref{prop:op-preclan-morphism-small} and \Cref{lem:vupp-eq}.
\end{proof}

\begin{axiom}\label{ax:vertical-opposite-small}
  The equivalent conditions of
  \Cref{prop:vertical-opposite-small} hold.
\end{axiom}

We can view the pullbacks of the vertical opposite of a general small
opfibration as follows.
\[ \begin{tikzcd}
	{\vOp[\Gamma]{(\Gamma \opext A)}} & {\Gamma \opext A^\op} & \Uupp & \Ulow \\
	& \Gamma & \U & \U
  \arrow["\iso"{description}, draw=none, from=1-1, to=1-2]
	\arrow[from=1-2, to=1-3]
	\arrow[from=1-2, to=2-2]
	\arrow["\lrcorner"{anchor=center, pos=0.125}, draw=none, from=1-2, to=2-3]
	\arrow[from=1-3, to=1-4]
	\arrow["\uupp"', from=1-3, to=2-3]
	\arrow["\lrcorner"{anchor=center, pos=0.125}, draw=none, from=1-3, to=2-4]
	\arrow["\ulow", from=1-4, to=2-4]
	\arrow["A", from=2-2, to=2-3]
	\arrow["{A^\op}"', bend right, from=2-2, to=2-4]
	\arrow["{\ulcorner \uupp \urcorner}", from=2-3, to=2-4]
\end{tikzcd} \]

We call $\vlow : \Vlow \to \V$ the \emph{universal small fibration}.
For a map $A : \Gamma \to \V$,
we have a chosen pullback of $\vlow : \Vlow \to \V$ along $A$
given by
\[\Gamma \ext A := \Op{(\Op{\Gamma} \opext (\ulcorner \uupp \urcorner \circ \Op{A}))}\]
\[ \begin{tikzcd}
    {\Op{\Gamma} \opext (\ulcorner \uupp \urcorner \circ \Op{A})} & \Uupp & \Ulow \\
    {\Op{\Gamma}} & \U & \U
    \arrow[from=1-1, to=1-2]
    \arrow[from=1-1, to=2-1]
    \arrow["\lrcorner"{anchor=center, pos=0.125}, draw=none, from=1-1, to=2-2]
    \arrow[from=1-2, to=1-3]
    \arrow["\uupp"', from=1-2, to=2-2]
    \arrow["\lrcorner"{anchor=center, pos=0.125}, draw=none, from=1-2, to=2-3]
    \arrow["\ulow", from=1-3, to=2-3]
    \arrow["{\Op{A}}"', from=2-1, to=2-2]
    \arrow["{\ulcorner \uupp \urcorner}"', from=2-2, to=2-3]
  \end{tikzcd}
  \quad \rightsquigarrow \quad
  \begin{tikzcd}
    {\Gamma \ext A} & \Vlow & \Vupp \\
    \Gamma & \V & \V
    \arrow[from=1-1, to=1-2]
    \arrow["{d_A}"', from=1-1, to=2-1]
    \arrow["\lrcorner"{anchor=center, pos=0.125}, draw=none, from=1-1, to=2-2]
    \arrow[from=1-2, to=1-3]
    \arrow["\vlow"', from=1-2, to=2-2]
    \arrow["\lrcorner"{anchor=center, pos=0.125}, draw=none, from=1-2, to=2-3]
    \arrow["\vupp", from=1-3, to=2-3]
    \arrow["A"', from=2-1, to=2-2]
    \arrow[from=2-2, to=2-3]
  \end{tikzcd}\]
We refer to $\Gamma \ext A$ as the \emph{fibrant context extension}.

% We can view the pullbacks of the vertical opposite of a general small
% (op)fibration as follows.
% \[ \begin{tikzcd}
% 	{\vOp[\Gamma]{(\Gamma \opext A)}} & {\Gamma \opext A^\op} & \Uupp & \Ulow \\
% 	& \Gamma & \U & \U
%   \arrow["\iso"{description}, draw=none, from=1-1, to=1-2]
% 	\arrow[from=1-2, to=1-3]
% 	\arrow[from=1-2, to=2-2]
% 	\arrow["\lrcorner"{anchor=center, pos=0.125}, draw=none, from=1-2, to=2-3]
% 	\arrow[from=1-3, to=1-4]
% 	\arrow["\uupp"', from=1-3, to=2-3]
% 	\arrow["\lrcorner"{anchor=center, pos=0.125}, draw=none, from=1-3, to=2-4]
% 	\arrow["\ulow", from=1-4, to=2-4]
% 	\arrow["A", from=2-2, to=2-3]
% 	\arrow["{A^\op}"', bend right, from=2-2, to=2-4]
% 	\arrow["{\ulcorner \uupp \urcorner}", from=2-3, to=2-4]
% \end{tikzcd} \]

% \[\begin{tikzcd}
% 	{\vOp[\Gamma]{(\Gamma \ext A)}} & {\Gamma \ext A^\op} & \Vupp & \Vlow \\
% 	& \Gamma & \V & \V
% 	\arrow["\iso"{description}, draw=none, from=1-1, to=1-2]
% 	\arrow[from=1-2, to=1-3]
% 	\arrow[from=1-2, to=2-2]
% 	\arrow["\lrcorner"{anchor=center, pos=0.125}, draw=none, from=1-2, to=2-3]
% 	\arrow[from=1-3, to=1-4]
% 	\arrow["\vupp"', from=1-3, to=2-3]
% 	\arrow["\lrcorner"{anchor=center, pos=0.125}, draw=none, from=1-3, to=2-4]
% 	\arrow["\vlow", from=1-4, to=2-4]
% 	\arrow["A", from=2-2, to=2-3]
% 	\arrow["{{A^\op}}"', bend right, from=2-2, to=2-4]
% 	\arrow["{{\ulcorner \vupp \urcorner}}", from=2-3, to=2-4]
% \end{tikzcd}
% \]

Suppose $A : X \to \Gamma$ in $\fib_\V$ is a small fibration.
Then $\Op{A} : \Op{X} \to \Op{\Gamma}$ is a small opfibration
and hence has a classifying map 
\[
\ulcorner \Op{A} \urcorner : \Op{\Gamma \to \U}
\]
We define the classifying map for $A$ by
\[ \ulcorner A \urcorner := \Op{(\ulcorner \uupp \urcorner \circ \ulcorner \Op{A} \urcorner)} : \Gamma \to \V\]
It follows that $A : X \to \Gamma$ is isomorphic to the extension
$\Gamma \ext \ulcorner A \urcorner \to \Gamma$ in $\fib(\Gamma)$.
\[ \begin{tikzcd}
    X & {\Gamma \ext \ulcorner A \urcorner} & \Vlow \\
    & \Gamma & \V
    \arrow["\sim", from=1-1, to=1-2]
    \arrow["A"', from=1-1, to=2-2]
    \arrow[from=1-2, to=1-3]
    \arrow[from=1-2, to=2-2]
    \arrow["\lrcorner"{anchor=center, pos=0.125}, draw=none, from=1-2, to=2-3]
    \arrow["\vlow", from=1-3, to=2-3]
    \arrow["{\ulcorner A \urcorner}"', from=2-2, to=2-3]
  \end{tikzcd} \]

\begin{definition}\label{def:ACM-universe}
  An algebraic categorical model (ACM)
  \[(\CC,\opfibcart,\Op,\vOp,\ulow)\]
  consists of an op structure (\Cref{def:ACM})
  and a universe $\ulow : \Ulow \to \U$ satisfying
  Axioms \labelcref{ax:universe,ax:classifying-map,ax:vertical-opposite-small}.
\end{definition}

\section*{$\Sigma$-types}
We fix an ACM $(\CC,\opfibcart,\Op,\vOp,\ulow)$.
Recall the definition of polynomial composition from \Cref{def:pcomp}.
We consider the polynomial composition of $(\ulow : \Ulow \to \U)$
in $\Poly_{(\CC,\fib)}$ with itself,
noting that $\U$ is fibrant, and $\ulow$ is $\fib$-exponentiable.
For notational convenience, we write $Q$ for the domain of the composition.
\[\begin{tikzcd}
	\Ulow & Q & \\
	\U & \fstProj \times_\U \ulow & \Ulow \\
	& {P_\ulow \U} & \U
	\arrow["\ulow"', from=1-1, to=2-1]
	\arrow[from=1-2, to=1-1]
	\arrow["\lrcorner"{anchor=center, pos=0.125, rotate=-90}, draw=none, from=1-2, to=2-1]
	\arrow[from=1-2, to=2-2]
	\arrow["{{\ulow \pcomp \ulow}}"{description, pos=0.2}, shift left=5, bend left, from=1-2, to=3-2]
	\arrow["{{{\sndProj}}}", from=2-2, to=2-1]
	\arrow[from=2-2, to=2-3]
	\arrow[from=2-2, to=3-2]
	\arrow["\lrcorner"{anchor=center, pos=0.125}, draw=none, from=2-2, to=3-3]
	\arrow["\ulow", from=2-3, to=3-3]
	\arrow["{{{\fstProj}}}"', from=3-2, to=3-3]
\end{tikzcd}\]
Similarly, we have the composition $\vlow \pcomp \vlow : Q' \to P_\vlow \V$.

\begin{theorem}\label{thm:sigma-equiv-defs}
  The following are equivalent (propositions)
  \begin{enumerate}
    \item The class of maps $\opfib_\U$ is closed under composition.
    \item $\ulow \pcomp \ulow : Q \to P_\ulow \U$ is a small opfibration.
    \item The class of maps $\fib_\V$ is closed under composition.
    \item $\vlow \pcomp \vlow : Q' \to P_\vlow \V$ is a small fibration.
  \end{enumerate}
  When these conditions hold,
  \Cref{ax:classifying-map} provides maps
  $\Sigma : P_\ulow \U \to \U$
  and $\pair : Q \to \Ulow$ forming a pullback square as follows (and dually).
  \begin{equation}\label{eq:fib-algebraic-sigma}
    \begin{tikzcd}
      Q & \Ulow \\
      {P_\ulow \U} & \U
      \arrow["{\pair}", from=1-1, to=1-2]
      \arrow["{{\ulow \pcomp \ulow}}"', from=1-1, to=2-1]
      \arrow["\lrcorner"{anchor=center, pos=0.125}, draw=none, from=1-1, to=2-2]
      \arrow["\ulow", from=1-2, to=2-2]
      \arrow["{\Sigma}"', from=2-1, to=2-2]
    \end{tikzcd}
    \quad 
    \begin{tikzcd}
      Q' & \Vlow \\
      {P_\vlow \V} & \V
      \arrow["\pair", from=1-1, to=1-2]
      \arrow["{{\vlow \pcomp \vlow}}"', from=1-1, to=2-1]
      \arrow["\lrcorner"{anchor=center, pos=0.125}, draw=none, from=1-1, to=2-2]
      \arrow["\vlow", from=1-2, to=2-2]
      \arrow["\Sigma"', from=2-1, to=2-2]
    \end{tikzcd}
  \end{equation}
\end{theorem}
\begin{proof}
  (1) $\Rightarrow$ (2) is immediate from the definition of $\ulow \pcomp \ulow$.
  (1) $\Leftrightarrow$ (3) is an immediate corollary of
  \Cref{prop:op-preclan-morphism-small}.
  (3) $\Leftrightarrow$ (4) is dual to (1) $\Leftrightarrow$ (2).
  Thus it suffices to prove (2) $\Rightarrow$ (1).
  
  If $d_C : C \to B$ and $d_B : B \to A$
  are two small opfibrations
  then \Cref{ax:classifying-map} provides classifying maps
  \[ \begin{tikzcd}
    \Ulow & C & \\
    \U & B & \Ulow \\
    & A & \U
    \arrow["\ulow"', from=1-1, to=2-1]
    \arrow["c"', from=1-2, to=1-1]
    \arrow["\lrcorner"{anchor=center, pos=0.125, rotate=-90}, draw=none, from=1-2, to=2-1]
    \arrow[from=1-2, to=2-2]
    \arrow["{\ulcorner C \urcorner}", from=2-2, to=2-1]
    \arrow["b", from=2-2, to=2-3]
    \arrow[from=2-2, to=3-2]
    \arrow["\lrcorner"{anchor=center, pos=0.125}, draw=none, from=2-2, to=3-3]
    \arrow["\ulow", from=2-3, to=3-3]
    \arrow["{\ulcorner B \urcorner}"', from=3-2, to=3-3]
  \end{tikzcd} \]
  By applying the universal property of polynomials from
  \Cref{prop:poly-up},
  we have that the composition $d_B \circ d_A : C \to A$ 
  is a pullback of $\ulow \pcomp \ulow$.
  Hence $d_B \circ d_A$ is also a small opfibration.
  \[\begin{tikzcd}
    C & Q \\
    B & {\fstProj \times_\U \ulow} \\
    A & {P_\ulow \U}
    \arrow["{(d_C, c)}", from=1-1, to=1-2]
    \arrow["{d_C}"', from=1-1, to=2-1]
    \arrow["\lrcorner"{anchor=center, pos=0.125}, draw=none, from=1-1, to=2-2]
    \arrow[from=1-2, to=2-2]
    \arrow["{(d_B, b)}", from=2-1, to=2-2]
    \arrow["{d_B}"', from=2-1, to=3-1]
    \arrow["\lrcorner"{anchor=center, pos=0.125}, draw=none, from=2-1, to=3-2]
    \arrow[from=2-2, to=3-2]
    \arrow["{({\ulcorner B \urcorner},{\ulcorner C \urcorner})}"', from=3-1, to=3-2]
  \end{tikzcd}\]
  % Thus $\ulow \pcomp \ulow : Q \to P_\ulow \U$ is classified by
  % $\ulow : \Ulow \to \U$.
  % \Cref{ax:classifying-map} provides maps $\Sigma : P_\ulow \U \to \U$
  % and $\pair : Q \to \Ulow$ forming a pullback square as follows.
\end{proof}
% \Cref{def:vert-cart-ofs} describes the $(\vert,\cart)$ factorisation system on
% $\Poly_{(\CC,\fib)}$.
Like in the theory of Martin-L\"of algebras \cite{awodey2025},
\cref{eq:fib-algebraic-sigma}
is a cartesian morphism of polynomial signatures in ${(\CC,\fib)}$.
In this sense,
the conditions of the theorem say that
the universe $\ulow : \Ulow \to \U$ has $\fib$-algebraic $\Sigma$-types
and the universe $\vlow : \Vlow \to \V$ has $\opfib$-algebraic $\Sigma$-types.

A similar analysis of $\Unit$-types can be given,
which we do not spell out.
We are led to the following axiom.

\begin{axiom}\label{ax:small-preclan}
  The class of maps $\opfib_\U$ is a preclan.
\end{axiom}
It follows that the equivalent conditions of \Cref{thm:sigma-equiv-defs}
hold and the class of maps $\fib_\V$ is a preclan.

In the $\Cat$ model (\Cref{ex:ACM-of-opfibrations}),
we will see that the $\Sigma$-type formation rule
is exactly the \emph{Grothendieck construction} operation.

\begin{definition}\label{def:ACM-universe-sigma}
  When an ACM also satisfies \Cref{ax:small-preclan},
  then we say that the universe admits $\Unit$-types and $\Sigma$-types.
\end{definition}

\section*{$\Pi$-types} 
% \label{sec:cat-types-pi-types}
We fix an ACM $(\CC,\opfibcart,\Op,\vOp,\ulow)$.
Consider the polynomial functor $P_\vlow : \opfib(1) \to \opfib(1)$
for the signature $(\vlow : \Vlow \to \V)$
in $\Poly_{(\CC,\opfib)}$.
We apply $P_\vlow$ to the morphism $\ulow : \Ulow \to \U$ in $\opfib(1)$.
\[ P_\vlow \ulow : P_\vlow \Ulow \to P_\vlow \U\]

\begin{theorem}\label{thm:pi-equiv-defs}
  The following are equivalent (propositions)
  \begin{enumerate}
    \item Small fibrations are small-opfibration-exponentiable. \[\fib_\V \subseteq \Exp_{\opfib_\U}\]
    \item $P_\vlow \ulow : P_\vlow \Ulow \to P_\vlow \U$ is a small opfibration.
    \item Small opfibrations are small-fibration-exponentiable. \[\opfib_\U \subseteq \Exp_{\fib_\V}\]
    \item $P_\ulow \vlow : P_\ulow \Vlow \to P_\ulow \V$ is a small fibration.
    \item $(\Cat,\opfib_\U, \fib_\V)$ is a $\tau$-clan.
  \end{enumerate}
  When these conditions hold,
  \Cref{ax:classifying-map} provides maps
  $\Pi : P_\vlow \U \to \U$ and $\lam : P_\vlow \Ulow \to \Ulow$
  forming the following pullback square (and dually).
\begin{equation}\label{eq:fib-algebraic-pi}
\begin{tikzcd}
	{P_\vlow \Ulow} & \Ulow \\
	{P_\vlow \U} & \U
	\arrow["\lam", from=1-1, to=1-2]
	\arrow["{P_\vlow \ulow}"', from=1-1, to=2-1]
	\arrow["\lrcorner"{anchor=center, pos=0.125}, draw=none, from=1-1, to=2-2]
	\arrow["\ulow", from=1-2, to=2-2]
	\arrow["\Pi"', from=2-1, to=2-2]
\end{tikzcd}
\quad \quad
\begin{tikzcd}
	{P_\ulow \Vlow} & \Vlow \\
	{P_\ulow \V} & \V
	\arrow["\lam", from=1-1, to=1-2]
	\arrow["{P_\ulow \vlow}"', from=1-1, to=2-1]
	\arrow["\lrcorner"{anchor=center, pos=0.125}, draw=none, from=1-1, to=2-2]
	\arrow["\vlow", from=1-2, to=2-2]
	\arrow["\Pi"', from=2-1, to=2-2]
\end{tikzcd}
\end{equation}
\end{theorem}
\begin{proof}
  (1) $\Leftrightarrow$ (3) $\Leftrightarrow$ (5) is the same as \Cref{prop:exponentiability-opposite}.
  (1) $\Leftrightarrow$ (2) and (3) $\Leftrightarrow$ (4) are dual,
  so we just show (1) $\Leftrightarrow$ (2).
  (1) $\Rightarrow$ (2) follows from viewing
  $(\vlow : \Vlow \to \V)$ as
  a signature in $\Poly_{(\CC,\opfib)}$
  and applying \Cref{lem:polynomial-functor-preserves-maps}
  with $\SS = \opfib_\U$ and $\RR = \opfib$.

  We sketch (2) $\Rightarrow$ (1).
  By \Cref{lem:stable-class-of-exponentiable-maps},
  it suffices to show that for a small opfibration $d_C : C \to B$ and
  a small fibration $d_B : B \to A$,
  there is a small opfibration $d_D : D \to A$
  satisfying the universal property of the pushforward of $d_C$ along $d_B$.
  Again, \Cref{ax:classifying-map} provides classifying maps
  \[\begin{tikzcd}
    C & \Vlow && B & \Ulow \\
    B & \V && A & \U
    \arrow[from=1-1, to=1-2]
    \arrow[from=1-1, to=2-1]
    \arrow["\lrcorner"{anchor=center, pos=0.125}, draw=none, from=1-1, to=2-2]
    \arrow["\vlow", from=1-2, to=2-2]
    \arrow[from=1-4, to=1-5]
    \arrow[from=1-4, to=2-4]
    \arrow["\lrcorner"{anchor=center, pos=0.125}, draw=none, from=1-4, to=2-5]
    \arrow["\ulow", from=1-5, to=2-5]
    \arrow["{{\ulcorner C \urcorner}}"', from=2-1, to=2-2]
    \arrow["{{\ulcorner B \urcorner}}"', from=2-4, to=2-5]
  \end{tikzcd}\]
  By applying the universal property of polynomials from \Cref{prop:poly-up},
  we have a map
  $(\ulcorner B \urcorner, \ulcorner C \urcorner) : A \to P_\vlow \U$.
  The map $d_D : D \to A$ is given by taking the pullback of
  $P_\vlow \ulow : P_\vlow \Ulow \to P_\vlow \U$ along this map.
  \[ \begin{tikzcd}[column sep = large]
      D & {P_\vlow \Ulow} \\
      A & {P_\vlow \U}
      \arrow[from=1-1, to=1-2]
      \arrow[from=1-1, to=2-1]
      \arrow["\lrcorner"{anchor=center, pos=0.125}, draw=none, from=1-1, to=2-2]
      \arrow["{P_\vlow \ulow}", from=1-2, to=2-2]
      \arrow["{(\ulcorner B \urcorner, \ulcorner C \urcorner)}"', from=2-1, to=2-2]
    \end{tikzcd} \]
\end{proof}

\begin{axiom}\label{ax:pi}
  The equivalent conditions of \Cref{thm:pi-equiv-defs} hold.
\end{axiom}

Since $P_\vlow \U$ is in $\opfib(1) = \fib(1)$,
the map $(P_\vlow \ulow : P_\vlow \Ulow \to P_\vlow \U)$ is also a
polynomial signature in $(\CC,\fib)$,
making \cref{eq:fib-algebraic-pi} a cartesian morphism of signatures.
Thus, the universe $\ulow : \Ulow \to \U$ has $\fib$-algebraic $\Pi$-types 
and the universe $\vlow : \Vlow \to \V$ has $\opfib$-algebraic $\Pi$-types.

\begin{definition}\label{def:ACM-universe-pi}
  When an ACM also satisfies \Cref{ax:pi},
  then we say that the universe admits $\Pi$-types.
\end{definition}

\section*{Hom-types}
% \label{sec:cat-types-hom-types}
We fix an ACM $(\CC,\opfibcart,\Op,\vOp,\ulow)$.
In this section,
we axiomatise a directed version of algebraic identity types
\cite{awodey2018, awodey2025},
also drawing from the analysis of algebraic path-types in
\cite{awodey2026}.
Before doing so,
we require additional external structure on $\CC$.

\begin{axiom} \label{ax:core}
  The pseudofunctor $\opfibcart(-) : \CC^\op \to \tCat$
  for opfibrations admits the following \emph{core} structure.
  \begin{itemize}
  \item We have a pseudotransformation
    \[\vCore{} : \opfibcart(-) \to \opfibcart(-)\]
  \item We have two modifications
    \[ \inv : \vCore{} \to \vOp{}
    \quad \text{and} \quad
    \inc : \vCore{} \to \id_{\opfibcart(-)} \]
  % \item The whiskerings
  % \[\inv \vCore{} : \vCore{} \vCore{} \to \vOp{} \vCore{}\]
  % and
  % \[\inc \vCore{} : \vCore{} \vCore{} \to \vCore{}\]
  % are isomorphisms.
  % We will write $\cml : \vCore{} \to \vCore{}\vCore{}$
  % for the inverse of $\inc \vCore{}$.
  \end{itemize}
\end{axiom}
For an opfibration $A \to X$,
the opfibration $\vCore[X]{A} \to X$
is the \emph{vertical core} of $A$ over $X$.
In $\Cat$, the vertical core of $A$ can be thought of as
the category where,
over each point $x$ in $X$,
the fibre is the groupoid core of $A_x$.

\begin{axiom} \label{ax:twist}
  The pseudofunctor $\opfibcart(-) : \CC^\op \to \tCat$
  for opfibrations admits the following \emph{twisted} structure.
  \begin{itemize}
  \item We have a pseudotransformation 
    \[\vTw{} : \opfibcart(-) \to \opfibcart(-)\]
  \item We have three modifications
    \[ \src : \vTw{} \to \vOp{}
      \quad \quad
      \trg : \vTw{} \to \id_{\opfibcart(-)}  \]
    \[ \rfl : \vCore{} \to \vTw{} \]
  \item These modifications satisfy two equations
    $\src \circ \rfl = \inv$ and $\trg \circ \rfl = \inc$.
    \[
      \begin{tikzcd}
        & {\vOp{}} \\
        {\vCore{}} & {\vTw{}} \\
        & {\id_{\opfibcart(-)}}
        \arrow["\inv", from=2-1, to=1-2]
        \arrow["\rfl"{description}, from=2-1, to=2-2]
        \arrow["\inc"', from=2-1, to=3-2]
        \arrow["\src"', from=2-2, to=1-2]
        \arrow["\trg", from=2-2, to=3-2]
      \end{tikzcd}
    \]
  \end{itemize}
  We call $\vTw[X]{A}$ the \emph{(vertical) twisted arrows} of $A$ over $X$.
\end{axiom}

For any object $X \in \CC$ and any opfibration $A \in \opfibcart(X)$ over it,
we have a \emph{twisted pathobject factorisation} of the (restricted) diagonal
$\Delta : \vCore[X]{A} \to \vOp[X]{A} \times_X A$.
\[\begin{tikzcd}
	{\vCore[X] A} & {\vTw{A}} \\
	& {\vOp[X]{A} \times_X A}
	\arrow["\rfl", from=1-1, to=1-2]
	\arrow["{\Delta := (\inv,\inc)}"', from=1-1, to=2-2]
	\arrow["{(\src,\trg)}", from=1-2, to=2-2]
\end{tikzcd}\]

\begin{axiom}\label{ax:src-trg-opfibration}
  For each $X \in \CC$ and $A \in \opfibcart(X)$,
  the map $(\src,\trg) : \vTw{A} \to \vOp[X]{A} \times_X A$
  is an opfibration.
\end{axiom}

Applying $\vCore[U]{}$ and $\vTw[U]{}$ to the universe of small opfibrations
$\ulow : \Ulow \to \U$, we have
\[ \begin{tikzcd}
	\Usim &&& \Utw \\
	\U &&& \U
	\arrow["{\usim \, := \, \vCore[U]{\ulow}}", from=1-1, to=2-1]
	\arrow["{\utw \, := \, \vTw[U]{\ulow}}", from=1-4, to=2-4]
\end{tikzcd} \]
% Then the modifications
% $\src : \Tw{} \to \vOp{}$ and $\trg : \Tw{} \to \id_{\opfibcart(-)}$
% By \Cref{ax:src-trg-opfibration}
% we have the following 
% universal small twisted arrow type.
We consider the universal small twisted pathobject factorisation.
\[\begin{tikzcd}
	\Usim & \Utw \\
	& {\Uupp \times_\U \Ulow}
	\arrow["\rfl", from=1-1, to=1-2]
	\arrow["\Delta"', from=1-1, to=2-2]
	\arrow["{{(\src,\trg)}}", from=1-2, to=2-2]
\end{tikzcd}\]

\begin{proposition}\label{prop:src-trg-small-opfibration}
  The following are equivalent
  \begin{itemize}
    \item For each $X \in \CC$ and small opfibration $A \in \opfibcart_\U(X)$,
      the map $(\src,\trg) : \vTw{A} \to \vOp[X]{A} \times_X A$
      is a small opfibration.
    \item $(\src,\trg) : \Utw \to \Uupp \times_\U \Ulow$
      is a small opfibration.
  \end{itemize}
  When these conditions hold,
  \Cref{ax:classifying-map} provides maps
  $\Dir : \Uupp \times_\U \Ulow \to \U$
  and $\dir : \Utw \to \Ulow$ such that
  the following square is a pullback.
  \[
    \begin{tikzcd}
      \Utw & \Ulow \\
      {\Uupp \times_\U \Ulow} & \U
      \arrow["\dir", from=1-1, to=1-2]
      \arrow["{(\src, \trg)}"', from=1-1, to=2-1]
      \arrow["\ulow", from=1-2, to=2-2]
      \arrow["\Dir"', from=2-1, to=2-2]
      \arrow["\lrcorner"{anchor=center, pos=0.125}, draw=none, from=1-1, to=2-2]
    \end{tikzcd}
  \]
\end{proposition}
\begin{proof}
  As usual, every operation in $\opfibcart_\U(X)$
  is a pullback of a universal one in $\opfibcart_\U(\U)$.
\end{proof}

This gives us a generalisation of the condition (A1) in \cite{awodey2026},
captured in the following axiom.
\begin{axiom}\label{ax:hom-formation-introduction}
  The equivalent conditions of
  \Cref{prop:src-trg-small-opfibration} hold.
\end{axiom}

From \Cref{ax:hom-formation-introduction},
we obtain formation and introduction rules for hom-types.
We now turn our attention to the hom-type elimination rule.

\begin{definition}
  We define the pseudotransformation
  $\rtw{} : \opfibcart(-) \to \opfibcart(-)$
  such that for each $X$ in $\CC$
  and object $A$ in $\opfibcart(X)$,
  $\rtw[X]{A}$ in $\opfibcart(X)$
  is the pullback in the functor category.
  Indeed, by \Cref{ax:src-trg-opfibration}
  the pullback
  $(\src,\trg) : \rtw[X]{A} \to \vCore[X]{A} \times_X A$
  exists and is an opfibration.
  \[
    \begin{tikzcd}
      {\rtw[X]{A}} & {\vTw[X]{A}} \\
      {\vCore[X]{A} \times_X A} & {\vOp[X]{A} \times_X A}
      \arrow[dashed, from=1-1, to=1-2]
      \arrow["{(\src,\trg)}"', dashed, from=1-1, to=2-1]
      \arrow["\lrcorner"{anchor=center, pos=0.125}, draw=none, from=1-1, to=2-2]
      \arrow["{(\src,\trg)}", from=1-2, to=2-2]
      \arrow["{\inv \times_X A}"', from=2-1, to=2-2]
    \end{tikzcd}
  \]
We call $\rtw[X] A$ the \emph{right-twisted arrows} of $A$ over $X$.
\end{definition}
We also obtain a \emph{right-twisted pathobject factorisation},
by the pullback property of $\rtw[X]{A}$.
\[\begin{tikzcd}
	{\vCore[X] A} & {\rtw[X]{A}} \\
	& {\vCore[X]{A} \times_X A}
	\arrow["\rfl", from=1-1, to=1-2]
	\arrow["{\Delta := (\id,\inc)}"', from=1-1, to=2-2]
	\arrow["{(\src,\trg)}", from=1-2, to=2-2]
\end{tikzcd}\]
There is, of course, also $\ltw[X] A$ the \emph{left-twisted arrows} of $A$ over $X$,
given by instead taking the pullback along the map
\[\vOp[X]{A} \times_A \inc : \vOp[X]{A} \times_X \vCore[X] A \to {\vOp[X]{A} \times_X A}\]

Applying $\rtw[\U]$ to $\ulow : \Ulow \to \U$
produces a universal small right-twisted arrow type
$\Urtw := \rtw[\U]{\Ulow}$ in $\opfibcart(\U)$
with $\urtw : \Urtw \to \U$
and an opfibration $(\src,\trg) : \Urtw \to \Usim \times_\U \Ulow$.
  \[ \begin{tikzcd}[column sep=huge]
    	\Usim & \Urtw & \Utw \\
      & {\Usim \times_\U \Ulow} & {\Uupp \times_\U \Ulow}
      \arrow["\rfl", from=1-1, to=1-2]
      \arrow["\Delta"', from=1-1, to=2-2]
      \arrow[from=1-2, to=1-3]
      \arrow["{{(\src,\trg)}}"{description}, from=1-2, to=2-2]
      \arrow["\lrcorner"{anchor=center, pos=0.125}, draw=none, from=1-2, to=2-3]
      \arrow["{{(\src,\trg)}}", from=1-3, to=2-3]
      \arrow["{{\inv \times_\U \Ulow}}"', from=2-2, to=2-3]
    \end{tikzcd} \]
Furthermore, $(\src,\trg) : \Urtw \to \Usim \times_\U \Ulow$
is a small opfibration by \Cref{ax:hom-formation-introduction}.
Consider the following commuting triangle.
\[ \begin{tikzcd}
	\Usim & \Urtw \\
	& \U
	\arrow["\rfl", from=1-1, to=1-2]
	\arrow["{\usim}"', from=1-1, to=2-2]
	\arrow["\urtw", from=1-2, to=2-2]
\end{tikzcd} \]
Both $(\usim : \Usim \to \U)$ and $(\urtw : \Urtw \to \U)$ are
polynomial signatures in $\Poly_{(\CC,\FF)}$,
since they are both opfibrations.
Using \Cref{def:vert-trans},
this triangle gives rise to a vertical natural transformation
between the polynomial functors.
\[P_{\rfl} : P_{\urtw} \to P_{\usim}\]
We consider the naturality square of $P_{\rfl}$
for the morphism $\ulow : \Ulow \to \U$.
\Cref{ax:hom-elimination}
will make this square a ``weak pullback''.
\[ \begin{tikzcd}
	{ P_{\urtw} \Ulow} & {P_{\usim} \Ulow} \\
	{ P_{\urtw} \U} & {P_{\usim} \U}
	\arrow["{{P_{\rfl} \Ulow}}", from=1-1, to=1-2]
	\arrow["{{ P_{\urtw} \ulow}}"', from=1-1, to=2-1]
	\arrow["{{P_{\usim} \ulow}}", from=1-2, to=2-2]
	\arrow["{{P_{\rfl} \U}}"', from=2-1, to=2-2]
\end{tikzcd} \]
We would like to compare $P_{\urtw}$ with an actual pullback
of the cospan in question.
Unfortunately, the other axioms do not automatically provide
such a pullback, so we postulate one.

\begin{axiom}\label{ax:pullback-L}
  There is a pullback $L$,
  with maps $\phi$ and $\psi$ as depicted in the following diagram.
  \begin{equation}
    \label{eq:hom-elimination-2}
    \begin{tikzcd}
    { P_{\urtw} \Ulow} && \\
    & L & {P_{\usim} \Ulow} \\
    & { P_{\urtw} \U} & {P_{\usim} \U}
    \arrow["p", dashed, from=1-1, to=2-2]
    \arrow["{P_{\rfl} \Ulow}", bend left, from=1-1, to=2-3]
    \arrow["{P_{\urtw} \ulow}"', bend right, from=1-1, to=3-2]
    \arrow["\phi", from=2-2, to=2-3]
    \arrow["\psi"', from=2-2, to=3-2]
    \arrow["\lrcorner"{anchor=center, pos=0.125}, draw=none, from=2-2, to=3-3]
    \arrow["{P_{\usim} \ulow}", from=2-3, to=3-3]
    \arrow["{P_{\rfl} \U}"', from=3-2, to=3-3]
  \end{tikzcd}
  \end{equation}
  We denote the comparison map $p : P_{\urtw} \Ulow \to L$.  
\end{axiom}
By the universal property of polynomials
and commutativity of the square,
$\psi$ and $\phi$ are equivalent to the data of three maps
\[\psi_1 = \phi_1 : L \to \U
\quad \psi_2 : \rtw[L]{A} \to \U
\quad \phi_2 : \vCore[L]{A} \to \Ulow\]
where $A := L \opext \psi_1 \to L$ is the opfibration over $L$
given by pullback of $\ulow : \Ulow \to \U$ along $\psi_1 : L \to \U$.
The latter two maps fit into a commutative square
\begin{equation}\label{eq:hom-elimination-1}
  \begin{tikzcd}
	\Ulow & {\vCore[L]{A}} & \Usim \\
	\U & {\rtw[L]{A}} & \Urtw \\
	& L & \U
	\arrow["\ulow"', from=1-1, to=2-1]
	\arrow["{\phi_2}"', from=1-2, to=1-1]
	\arrow[from=1-2, to=1-3]
	\arrow["{\rfl}", from=1-2, to=2-2]
	\arrow["\lrcorner"{anchor=center, pos=0.125}, draw=none, from=1-2, to=2-3]
	\arrow["\rfl", from=1-3, to=2-3]
	\arrow["{\usim}", bend left = 50, from=1-3, to=3-3]
	\arrow["{\psi_2}", from=2-2, to=2-1]
	\arrow[from=2-2, to=2-3]
	\arrow[from=2-2, to=3-2]
	\arrow["\lrcorner"{anchor=center, pos=0.125}, draw=none, from=2-2, to=3-3]
	\arrow["\urtw", from=2-3, to=3-3]
	\arrow["{\psi_1}"', from=3-2, to=3-3]
\end{tikzcd}
\end{equation}

\begin{proposition}\label{prop:hom-elimination-equiv}
  The following are equivalent (structures)
  \begin{itemize}
    \item A section $j : L \to P_{\urtw} \Ulow$
    of the comparison map $p : P_{\urtw} \Ulow \to L$,
    from \cref{eq:hom-elimination-2}.
    \[ p \circ j = \id\]
    This is akin to \cite[eq. (12)]{awodey2025}
    used in the elimination principle for algebraic identity types.
    \item A diagonal lift $j : \rtw[L]{A} \to \Ulow$
    as indicated in the following diagram,
    extracted from \cref{eq:hom-elimination-1}.
    \[\begin{tikzcd}
        {\vCore[L]{A}} & \Ulow \\
        {\rtw[L]{A}} & \U
        \arrow["{{\phi_2}}", from=1-1, to=1-2]
        \arrow["{{\rfl}}"', from=1-1, to=2-1]
        \arrow["\ulow", from=1-2, to=2-2]
        \arrow["j"{description}, dashed, from=2-1, to=1-2]
        \arrow["{{\psi_2}}"', from=2-1, to=2-2]
      \end{tikzcd}\]
    \item A diagonal lift $j' : \rtw[L]{A} \to C$
    as indicated in the following diagram.
    \[\begin{tikzcd}
        {\vCore[L]{A}} & C & \Ulow \\
        {\rtw[L]{A}} & {\rtw[L]{A}} & \U
        \arrow["r", dashed, from=1-1, to=1-2]
        \arrow["{{{\phi_2}}}", bend left, from=1-1, to=1-3]
        \arrow["{{{\rfl}}}"', from=1-1, to=2-1]
        \arrow[from=1-2, to=1-3]
        \arrow[from=1-2, to=2-2]
        \arrow["\lrcorner"{anchor=center, pos=0.125}, draw=none, from=1-2, to=2-3]
        \arrow["\ulow", from=1-3, to=2-3]
        \arrow["{j'}"{description}, dashed, from=2-1, to=1-2]
        \arrow[equals, from=2-1, to=2-2]
        \arrow["{{{\psi_2}}}"', from=2-2, to=2-3]
      \end{tikzcd}\]
  \end{itemize}
\end{proposition}
\begin{proof}
  (1) $\Leftrightarrow$ (2) follows from the universal property of polynomial functors.
  (2) $\Leftrightarrow$ (3) is clear.
\end{proof}

The equivalent conditions of \Cref{prop:hom-elimination-equiv}
can be presented using an informal syntax as follows.

\begin{prooftree}
% \AxiomC{$\Gamma \vdash A : U$}
\AxiomC{$\Gamma \opext x : \vCore A \opext y : A \opext h : \Hom_A (\inv x, y) \vdash C(x,y,h) : U$}
\noLine
\UnaryInfC{$\Gamma \opext x : \vCore A \vdash r : C(x,\inc x, \rfl x)$}
\UnaryInfC{$\Gamma \opext x : \vCore A \opext y : A \opext h : \Hom_A (\inv x, y) \vdash j_{(C,r)}(x,y,h) : C$}
\noLine
\UnaryInfC{$\Gamma \opext x : \vCore A \vdash j_{(C,r)}(x,\inc x, \rfl x) \equiv r : C(x,\inc x, \rfl x)$}
\end{prooftree}

This is essentially the syntax from \cite{north2019},
slightly simplified due to the presence of $\Pi$-types.
Rather than providing a formal syntax,
we will reformulate \Cref{prop:hom-elimination-equiv}
in ``unalgebraic'' terms in \Cref{def:elementary-hom-type},
as was done in \Cref{sec:hottlean}.

\begin{axiom}\label{ax:hom-elimination}
  We have one of the equivalent structures from \Cref{prop:hom-elimination-equiv}.
  Also, we have the analogous structure for the left-twisted arrows $\ltw$,
  which we do not spell out.
\end{axiom}

% Like in \cite{hua2026},
% \Cref{ax:hom-elimination} could be unfolded into an ``elementary'' condition instead:
% for all objects $\Gamma$, all maps $A : \Gamma \to \U$,
% all maps $C : \Gamma \opext x : \vCore A \opext y : A \opext \Hom_A (\inv x, y) \to \U$
% and maps $r : \Gamma \opext x : \vCore A \to \Ulow$
% such that $\ulow \circ r = C \circ RRRRR$,
% there is a map $j : \Gamma \opext x : \vCore A \opext y : A \opext \Hom_A (\inv x, y) \to \Ulow$
% satisfying a typing rule $\ulow \circ j = C$,
% a computation rule $j \circ RRRRR = r$,
% as well as a substitution stability condition.

\begin{definition}\label{def:ACM-universe-hom}
  When an ACM also satisfies Axioms
  \labelcref{ax:core,,ax:twist,,ax:src-trg-opfibration,,ax:hom-formation-introduction,ax:pullback-L,ax:hom-elimination},
  then we say that the universe admits $\Hom$-types.
\end{definition}

\section{Unalgebraic categorical models}

% Rather than presenting a syntax to interpret into an ACM,
% which comes with its own set of difficulties,
% we provide a more syntax reformulation of an ACM,
% called an elementary categorical model (UCM) (following \cite{hua2026}).
% The idea is that when constructing a model,
% one would construct it as an ACM,
% whereas one would target an UCM when interpreting syntax into a model.
Rather than presenting a syntax to interpret into an ACM,
we provide an ``unalgebraic'' reformulation of the definitions in algebraic categorical models
that is more syntactic in style.
Like in \Cref{sec:hottlean},
we will introduce ``unalgebraic categorical models''.
The idea is that one would typically construct models as ACMs,
as ACMs are closer to how they appear ``in the wild''.
At the same time, defining an interpretation of the syntax into the model
is straightforward using an UCM.
An appropriate translation from one to the other is therefore useful.
This approach also avoids getting into the intricacies of designing a well-behaved syntax
for dependent type theory,
which is a challenge in and of itself.

\begin{definition}
  An unalgebraic categorical model (UCM) 
  \[(\CC,\U,\V,\ulow,\uupp,\usim,\vlow,\vupp,\vsim,\inc,\inv)\]
  consists of the following.
  \begin{itemize}
  \item A category $\CC$ with a terminal object.
  \item Three universes $\ulow : \Ulow \to \U, \uupp:\Uupp \to \U$ and $\usim : \Usim \to \U$ in $\CC$,
    in the sense of \Cref{def:elementary-universe}.
  \item Three universes $\vlow : \Vlow \to \V, \vupp:\Vupp \to \V$ and $\vsim : \Vsim \to \V$ in $\CC$.
  \item Morphisms $\inv : \usim \to \uupp$ and $\inc : \usim \to \ulow$ over $\U$, and
    $\inv : \vsim \to \vupp$ and $\inc : \vsim \to \vlow$ over $\V$.
    \[
    \begin{tikzcd}
      \Uupp & \Usim & \Ulow \\
      & \U
      \arrow["\uupp"', from=1-1, to=2-2]
      \arrow["\inv"', from=1-2, to=1-1]
      \arrow["\inc", from=1-2, to=1-3]
      \arrow["\usim"{description}, from=1-2, to=2-2]
      \arrow["\ulow", from=1-3, to=2-2]
    \end{tikzcd}
    \quad \quad 
    \begin{tikzcd}
      \Vupp & \Vsim & \Vlow \\
      & \V
      \arrow["\vupp"', from=1-1, to=2-2]
      \arrow["\inv"', from=1-2, to=1-1]
      \arrow["\inc", from=1-2, to=1-3]
      \arrow["\vsim"{description}, from=1-2, to=2-2]
      \arrow["\vlow", from=1-3, to=2-2]
    \end{tikzcd}
    \]
  \end{itemize}
\end{definition}
Like in \Cref{sec:UnstrSem},
an object of $\CC$ will be called a \emph{context} and a morphism will be called a substitution.
A map $A : \Gamma \to \U$ will be called a \emph{covariant type},
and a map $A : \Gamma \to \V$ will be called a \emph{contravariant type}.
Each universe admits its own kind of context extension.
For a covariant type $A : \Gamma \to \U$,
let us introduce the suggestive notation for context extensions
$\Gamma \opext (x : A)$ and $\Gamma \opext (x : \vOp A)$, and $\Gamma \opext (x : \vCore A)$
respectively for $\ulow, \uupp$ and $\usim$.
Similarly, we denote the context extensions for $\vlow, \vupp, \vsim$
respectively as $\Gamma \ext A$, $\Gamma \ext \vOp A$, and $\Gamma \ext \vCore A$
(or with explicit variables).

\begin{definition}[Unalgebraic $\Hom$-types] \label{def:elementary-hom-type}
  Let $(\CC,\U,\V,\dots)$ be an UCM.
  Then an \emph{unalgebraic $\Hom$-type structure}
  on $\U$ consists of the following.
  \begin{enumerate}
  \item
    For any context $\Gamma$, covariant type $A : \Gamma \to \U$,
    maps $a_0 : \Gamma \to \Uupp$ and
    $a_1 : \Gamma \to \Ulow$ such that
    $\uupp \circ a_0 = A$ and $\ulow \circ a_1 = A$,
    there is a type $\Hom_A (a_0,a_1) : \Gamma \to \U$.
  \item
    The construction $\Hom_A (a_0,a_1)$ is stable under substitution,
    meaning that for any $\sigma : \Delta \to \Gamma$ we have
    \[ \Hom_{A \circ \sigma} (\sigma \gg a_0, \sigma \gg a_1)
    = \sigma \gg \Hom_A (a_0,a_1) \]
  \item
    For any context $\Gamma$, covariant type $A : \Gamma \to \U$,
    and map $a : \Gamma \to \Usim$ such that
    $\usim \circ a = A$,
    there is a map $\rfl a : \Gamma \to \Ulow$
    satisfying
    $\ulow \circ \rfl a  = \Hom_A (\inv \circ \, a,\inc \circ \, a)$.
  \item The construction $\rfl$ is stable under substitution,
    meaning that for any $\sigma : \Delta \to \Gamma$ we have
    \[ \rfl (\sigma \circ a) = \sigma \circ \rfl a \]
  \item
    Given a covariant type $A : \Gamma \to \U$,
    one can construct the context 
    \[\Gamma \opext \vCore A \opext A \opext \Hom_A := \Gamma \opext (x : \vCore A) \opext (y : A) \opext \Hom_A(\inv \circ \, x,y)\]
    and the substitution
    \[\rho : \Gamma \opext (x : \vCore A) \to \Gamma \opext (x : \vCore A) \opext (y : A) \opext \Hom_A(\inv \circ \, x,y)\]
    \[\rho := \id_{(\Gamma \opext \vCore A)} \opext (\inc \circ \, x) \opext \rfl x\]
    For any \emph{motive} \[C : \Gamma \opext \vCore A \opext A \opext \Hom_A \to \U\]
    and map $r : \Gamma \opext \vCore A \to \Ulow$ satisfying
    \[ \ulow \circ r = C \circ \rho \]
    there is a map, called the \emph{right eliminator},
    $j (C, r) : \Gamma \opext \vCore A \opext A \opext \Hom_A \to \Ulow$
    satisfying $\ulow \circ j(C,r) = C$ and
    $j (C, r) \circ \rho  = r$
    \[\begin{tikzcd}
	   \Gamma \opext \vCore A & \Ulow \\
	   {\Gamma \opext \vCore A \opext A \opext \Hom_A} & \U
	   \arrow["{r}", from=1-1, to=1-2]
	   \arrow["{\rho}"', from=1-1, to=2-1]
	   \arrow["\ulow", from=1-2, to=2-2]
	   \arrow["{j(C,r)}"{description}, dashed, from=2-1, to=1-2]
	   \arrow["C"', from=2-1, to=2-2]
    \end{tikzcd}\]
  \item The construction $j$ is stable under substitution.
  \item
    Dually, given a covariant type $A : \Gamma \to \U$,
    one can construct the context 
    \[\Gamma \opext \vOp A \opext \vCore A \opext \Hom_A :=
    \Gamma \opext (x : \vOp A) \opext (y : \vCore A) \opext \Hom_A(x,\inc \circ \, y)\]
    and the substitution
    \[\rho' : \Gamma \opext (x : \vCore A) \to \Gamma \opext \vOp A \opext \vCore A \opext \Hom_A\]
    \[\rho' := \id_{\Gamma} \opext (\inv \circ \, x) \opext x \opext \rfl x\]
    For any \emph{motive} \[C : \Gamma \opext \vOp A \opext \vCore A \opext \Hom_A \to \U\]
    and map $r : \Gamma \opext \vCore A \to \Ulow$ satisfying
    \[ \ulow \circ r = C \circ \rho' \]
    there is a map, called the \emph{left eliminator},
    $j' (C, r) : \Gamma \opext \vOp A \opext \vCore A \opext \Hom_A \to \Ulow$
    satisfying $\ulow \circ j'(C,r) = C$ and
    $j' (C, r) \circ \rho'  = r$
    \[\begin{tikzcd}
	   \Gamma \opext \vCore A & \Ulow \\
	   {\Gamma \opext \vOp A \opext \vCore A \opext \Hom_A} & \U
	   \arrow["{r}", from=1-1, to=1-2]
	   \arrow["{\rho'}"', from=1-1, to=2-1]
	   \arrow["\ulow", from=1-2, to=2-2]
	   \arrow["{j'(C,r)}"{description}, dashed, from=2-1, to=1-2]
	   \arrow["C"', from=2-1, to=2-2]
    \end{tikzcd}\]
  \item The construction $j'$ is stable under substitution.
  \end{enumerate}
\end{definition}

\begin{theorem}
  Given an algebraic categorical model with a universe that has
  $\Unit$-types, $\Sigma$-types, $\Pi$-types, and $\Hom$-types,
  there is an unalgebraic categorical model $(\CC,\U,\V,\dots)$
  with a $\vlow$-indexed {unalgebraic $\Pi$-type structure} on $\ulow$ (\Cref{def:elementary-pi}),
  a $\ulow$-indexed {unalgebraic $\Pi$-type structure} on $\vlow$,
  an {unalgebraic $\Sigma$-type structure} on $\ulow$ (\Cref{def:elementary-sigma}),
  an {unalgebraic $\Sigma$-type structure} on $\vlow$,
  an {unalgebraic $\Hom$-type structure} on $\U$,
  and an {unalgebraic $\Hom$-type structure} on $\V$.
\end{theorem}
\begin{proof}
  The proof of this is routine;
  similar proofs can be found in \cite{awodey2018,awodey2025,hua2026}.
\end{proof}

\section{Examples}

Models of Martin-L\"of type theory are degenerate categorical algebraic models.
\begin{example}[Models of Martin-L\"of type theory]
  Suppose $\CC$ is an algebraic model of Martin-L\"of type theory
  on a $\pi$-clan $(\CC,\RR)$.
  Let us take $\opfib := \RR$,
  and $\Op : \CC \to \CC$ to be the identity,
  so that we have $\RR = \opfib = \fib$.
  The $\tau$-clan condition in \Cref{ax:exponentiability}
  follows from $\RR$ being a $\pi$-preclan.
  We also set $\opfibcart := \RR$ so that all maps in the slice are
  morphisms of opfibrations.
  We then take $\vOp : \RR(-) \to \RR(-)$ to be the identity map.
  This results in an op structure, which we will call $\MM_\RR$.
  
  It follows that an ACM extending this op structure (\Cref{def:ACM-universe}) is the same thing as
  an ``$\RR$-algebraic universe'' in an algebraic model of MLTT (\Cref{def:algebraic-universe}).
  Let us assume now that we have such a universe $\uu : \Ulow \to \U$ in $\MM_\RR$.
  Then $\uu : \Ulow \to \U$ admits $\Unit$-types and $\Sigma$-types as an ACM (\Cref{def:ACM-universe-sigma})
  if and only if it does so in the MLTT model (\Cref{def:alg-unit-type-mltt,def:alg-sig-type-mltt}),
  and $\uu : \Ulow \to \U$ admits $\Pi$-types as an ACM (\Cref{def:ACM-universe-pi})
  if and only if it does so in the MLTT model (\Cref{def:alg-pi-type-mltt}).

  For $\Hom$-types,
  let us suppose that the universe in the MLTT model admits $\Path$-types
  (in the sense of \Cref{def:mltt-alg-path-types-2}).
  Take $\vCore : \RR(-) \to \RR(-)$ to be identity
  and the twisted arrows map $\vTw : \RR(-) \to \RR(-)$
  to be the path functor $\arr : \RR(-) \to \RR(-)$ (\Cref{def:path-functor}).
  Assuming that the universe admits a normal Hurewicz structure (\Cref{def:hurewicz-defs}),
  then the universe admits $\Id$-types (\Cref{thm:id-elim}),
  which coincide with $\Hom$-types.
\end{example}

There are several different algebraic categorical models in $\Cat$ with 
$\Unit$-types, $\Sigma$-types, $\Pi$-types, and $\Hom$-types.
One is with covariant types as split opfibrations.
In the following,
we will present this model in a way that uses the axiom of choice,
namely to satisfy \Cref{ax:classifying-map},
the axiom for having classifying maps.
It should be possible to reformulate the general theory of algebraic models
in such a way that avoids choice.
However, we will leave these constructive considerations for future work.

\begin{example}[$\Cat$ and the universe of opfibrations]\label{ex:ACM-of-opfibrations}
  \Cref{ax:opfibration-preclan}.
  Disregarding size issues that can be resolved by postulating
  more Grothendieck universes,
  we can take $\CC$ to be the 1-category of large categories $\Cat$.
  We take $\opfib$ to be the class of split opfibrations (\Cref{def:split-opfibration}).
  These together form a preclan (\Cref{prop:preclans-in-cat}).

  \Cref{ax:global-op} ($\Op$).
  Of course, taking opposite categories is an endofunctor on $\Cat$.
  The class of maps $\fib$ defined above is exactly the class of split fibrations.

  \Cref{ax:exponentiability}.
  We showed that $\Cat$ with opfibrations and fibrations forms a $\tau$-clan in \Cref{prop:pushforward-in-cat}.

  \Cref{ax:opcartesian-map} (morphisms of opfibrations).
  For each $X$, $\opfibcart(X)$ has morphisms as morphisms of
  split opfibrations over $\CC$ \Cref{def:morphism-split-opfibration}.
  For any category $\CC$, the map $\CC \to 1$ is a split opfibration,
  and the lift of the identity in $1$ is unique -
  since split opfibrations are also {normal}.
  It follows that any functor is also a
  split opcartesian map over $1$.

  \Cref{ax:local-op} ($\vOp$).
  Adjusting \Cref{thm:opfibration-main} for size,
  we have for any category $X$,
  \[\opfibcart(X) \simeq [X,\Cat]\]
  Postcomposing with the endofunctor $\Op : \Cat \to \Cat$
  yields a functor
  \[ \vOp[X] : \opfibcart(X) \to \opfibcart(X) \]
  as described in \Cref{def:relativised-constructions},
  which was for \emph{small} opfibrations.

  {\bf So far, we have constructed an op structure on $\Cat$.
  Let us call this op structure $\MM_\Cat$.}

  \Cref{ax:universe} (universe of small opfibrations).
  The universal small split opfibration in $\Cat$
  \[\ulow : \pcat \to \cat\]
  is an opfibration.
  This is defined in \Cref{def:universal-small-split-opfibration}.
  All objects are fibrant.

  \Cref{ax:classifying-map} (classifying maps).
  This follows from the axiom of choice.
  % Every small split opfibration $A \to X$ in $\Cat$ has a classifying map
  % $A^* : X \to \cat$, which takes a point in $X$ to its fibre, which is a category.
  % The mapping
  % \[A \mapsto A^* : \smlopfibcart(X) \to [X,\cat]\]
  % is one of the functors making up the equivalence
  % described in \Cref{thm:opfibration-main}.
  % \[[X,\cat] \simeq \smlopfibcart(X)\]

  \Cref{ax:vertical-opposite-small}.
  The first condition of \cref{prop:vertical-opposite-small} follows from
  the fact that straightening and unstraightening preserve size.
  \[
  \begin{tikzcd}
    {[A,\cat]} & {\smlopfibcart(A)} \\
    {[A,\Cat]} & {\opfibcart(A)}
    \arrow["\simeq"{description}, draw=none, from=1-1, to=1-2]
    \arrow[hook, from=1-1, to=2-1]
    \arrow[hook, from=1-2, to=2-2]
    \arrow["\simeq"{description}, draw=none, from=2-1, to=2-2]
  \end{tikzcd} 
  \]
  {\bf Thus $\MM_\Cat$ and $\ulow : \pcat \to \cat$ together form an ACM}.
  
  \Cref{ax:small-preclan} ($\Unit$-types and $\Sigma$-types).
  $\Cat$ with the class of small split opfibrations
  $(\Cat,\smlopfib)$
  is a preclan (\Cref{prop:preclans-in-cat}).
  In this case, given a small category $A$ and a functor $B : A \to \cat$,
  the point $(A,B)$ in $P_{\ulow} \U$ is taken by
  $\Sigma : P_{\ulow} \U \to \U$
  to the \emph{Grothendieck construction} $\int_A B$ in $\U$.

  \Cref{ax:pi} ($\Pi$-types).
  We showed that $\fib_V = \smlfib \subseteq \Exp_{\opfib_\U} = \Exp_{\smlopfib}$
  in \Cref{prop:pushforward-in-cat}.  
 
  \Cref{ax:core} ($\vCore$) and \Cref{ax:twist} ($\vTw$).
  The relativised twisted arrow structure for small split opfibrations
  is described in \Cref{def:relativised-constructions}.
  The same construction works for large split opfibrations.

  \Cref{ax:src-trg-opfibration} and \Cref{ax:hom-formation-introduction}
  ($\Hom$-type formation and introduction).
  We verified in \Cref{prop:local-src-trg-discrete-opfibration}
  that when $A \to X$ is a split opfibration,
  $\vTw[X] A \to \vOp[X] A \times_X A$ is a discrete opfibration,
  and therefore a split opfibration,
  and that this preserves size.

  \Cref{ax:pullback-L}.
  $\Cat$ has all 1-limits, giving us the required pullback $L$.

  \Cref{ax:hom-elimination} ($\Hom$ elimination).
  It is easiest to check the second condition in
  \Cref{prop:hom-elimination-equiv},
  namely that we have a lift $j : {\vCore[L]{A}} \to \Ulow$
  for the square
  \[\begin{tikzcd}
      {\vCore[L]{A}} & \Ulow \\
      {\rtw[L]{A}} & \U
      \arrow["{{\phi_2}}", from=1-1, to=1-2]
      \arrow["{{\rfl}}"', from=1-1, to=2-1]
      \arrow["\ulow", from=1-2, to=2-2]
      \arrow["{{\psi_2}}"', from=2-1, to=2-2]
    \end{tikzcd}\]
  Recall from \Cref{thm:lari-opfib-awfs} that
  we have an AWFS where the
  right maps are split opfibrations, and the left maps are lari.
  Of course, $\ulow : \Ulow \to \U$ is a split opfibration.
  By \Cref{prop:rtw-equiv-K},
  ${\rfl} : {\vCore[L]{A}} \to {\rtw[L]{A}}$ is lari.
  Hence, the AWFS provides a lift $j : {\vCore[L]{A}} \to \Ulow$, as required.

  {\bf In summary, the op structure $\MM_\Cat$ has a universe --
  the universe of split opfibrations -- with
  $\Unit$-types, $\Sigma$-types, $\Pi$-types, and $\Hom$-types.}
\end{example}

Next, we continue working with the op structure $\MM_\Cat$,
but we switch the universe of opfibrations
for the universe of \emph{groupoidal opfibrations}.
% the third has covariant types as discrete opfibrations

\begin{example}[$\Cat$ and the universe of groupoidal opfibrations]\label{ex:groupoidal-opfibrations}
  A split groupoidal opfibration is a
  split opfibration with fibres that are groupoids.
  Taking the pullback of $\ulow : \pcat \to \cat$ along the inclusion
  of the category of small groupoids
  $\grpd \to \cat$ produces the universes
  of \emph{small groupoidal opfibrations}.
  \[\quad \quad \pgrpd \to \grpd\]
  Although (large) groupoidal opfibrations also form preclans,
  these preclans are not suitable candidates for $\opfib$,
  as $\grpd$ is not itself a groupoid (so it is not in $\opfibcart(1)$).
  Hence, we still use the (large non-groupoidal)
  opfibrations as the ambient preclan $\opfib$.

  {\bf The universe of small split groupoidal opfibrations in the op structure $\MM_\Cat$ admits
  $\Unit$-types, $\Sigma$-types, $\Pi$-types, and $\Hom$-types.}
  For $\Hom$-types, the $\vOp[]{}$ and $\vCore[]{}$ operations can be taken to be the identity.
  \Cref{ax:hom-formation-introduction} holds for both of these universes;
  as noted above $\vTw[X] A \to \vOp[X] A \times_X A$
  is always a discrete opfibration.
  \Cref{ax:hom-elimination} for the universes of groupoidal opfibrations
  is a corollary of the case for general opfibrations.
  Moreover, as we will see in \Cref{prop:lari-orthogonal-groupoidal-opfibration},
  the diagonal filler $j$ in \Cref{prop:hom-elimination-equiv}
  is unique up to unique isomorphism.
\end{example}

\begin{example}[$\Cat$ and the universe of discrete opfibrations]\label{ex:discrete-opfibrations}
  Recall that a discrete opfibration is a
  split opfibration with fibres that are sets.
  Taking the pullback of $\ulow : \pcat \to \cat$ along the inclusion
  of the categories of small sets
  $\set \to \cat$ produces the universe
  of \emph{small discrete opfibrations}.
  \[\pset \to \set\]
  One can also check that the universe of small discrete opfibrations (in the op structure $\MM_\Cat$) admits
  $\Unit$-types, $\Sigma$-types, $\Pi$-types, and $\Hom$-types.
\end{example}

  % It is also straightforward to check that
  % $\Cat$ with the class of small groupoidal opfibrations
  % $(\Cat,\opfib_{\grpd})$
  % and the class of small discrete opfibrations
  % $(\Cat,\opfib_{\grpd})$
  % are preclans.
  % It is also straightforward to check that
  % discrete opfibrations are preserved under pushforward along
  % discrete fibrations
  % and groupoidal opfibrations are preserved under pushforward along
  % groupoidal fibrations.
  % \[ \fib_{\set} \subseteq \Exp_{\opfib_{\set}}
  %   \quad \quad \text{and} \quad \quad
  %   \fib_{\grpd} \subseteq \Exp_{\opfib_{\grpd}}\]